% !TeX root = ../papers/nc-p3.tex
\section{楕円五次点グラフの三割線と射影中心}
\label{sec:elliptic-quintic-trisecants}
以下、四辺で左右の射影ファイバーが空か一点の標準有限Frobenius候補で、
点グラフが滑らかな楕円五次曲線 $E$ の場合を扱う。$j_i:E\to\mathbb P(V_i^*)$、
$V_i\subset H^0(L_i)$ は非常に豊富な四次元部分系である。
完全な五切断系からの射影中心 $o_i$ は、接線を含む割線多様体の外にある。

\begin{proposition}\label{prop:quintic-mixed-kernel}
$R_i$ は常に $V_i\otimes V_{i+1}\to H^0(L_iL_{i+1})$ の全核である。
\end{proposition}
\begin{proof}
$L=L_i,M=L_{i+1}$ とする。積の像を消す非零な汎関数 $\ell$ があれば、
完全切断系上の双線型形式 $\ell(st)$ の階数は高々2である。
実際、二つの超平面の積上で零なので行と列を一本ずつ除けば零行列になる。
Serre双対でこの汎関数に対応する拡大を
$0\to\mathcal O_E\to F\to LM\to0$ とする。$L^{-1}$ でひねった接続射は上の双線型形式であり、
$F L^{-1}$ は次数0、階数2で三つ以上の独立な切断を持つ。
その切断の作る部分束が一般点で階数2なら、独立な二切断の外積は次数零で非零、
ゆえにどこでも消えず、この束は自明で切断は二次元しかない。
よって切断は一つの飽和直線束に入り、その次数は少なくとも3である。
元に戻すと $F$ には次数8以上の直線束が入り、$LM$ への像は $LM(-Z)$、$\ell(Z)\le2$
（ここで長さを表す）となる。引戻し拡大は分裂するので、汎関数は $H^0(LM(-Z))$ を消す。
両 $V_i,V_{i+1}$ は $Z$ 上の全切断を与え、その積は $LM|_Z$ の全切断を張るから矛盾する。
従って積は全射で、核次元 $16-10=6$ により主張を得る。
\end{proof}

\begin{lemma}\label{lem:elliptic-quintic-trisecants}
楕円五次埋込み $j:E\hookrightarrow\mathbb P^3$ で、各 $t\in\operatorname{Pic}^3 E$ に
その類の三割線因子 $T_t$ がスキームとして一意に定まる。
$Z=\{(t,q):q\in T_t\}$ の両射影は有限平坦で、次数は順に3と2である。
\end{lemma}
\begin{proof}
$N=t$、$M=L/N$ と書く。非零切断 $s\in H^0(N)$ の因子が直線に入る条件は
$sH^0(M)\subset V$ である。$H^0(N)\to H^0(M)^*$ を射影中心 $o$ と乗法で定めると階数は2である。
階数が落ちれば、ある非零 $m\in H^0(M)$ で $o$ は $mH^0(N)$ を消し、
長さ2の $(m)_0$ の張る割線上に入るからである。従って核はただ一本である。
次数4の因子が空間直線に乗ることもない。$L$ からそれを引けば次数1で、二つの独立平面を含めない。

$E_t\times E_q$ 上で $N$ を正規化Poincaré束に第二因子の次数3の束を掛けて取り、
$M=\operatorname{pr}_q^*L/N$ とする。$N,M$ の第一射影での順像は階数3,2で次数は共に $-1$。
実際、高次順像が零で、Poincaré因子 $\Delta-E\times\{0\}-\{0\}\times E$ の自乗は $-2$、
両座標ファイバーとの交叉は零なので曲面Riemann--Rochで各Euler数が $-1$ になる。
上の全射の核は次数 $-2$ の束 $B$ で、$Z$ の束は $N\otimes\operatorname{pr}_t^* B^{-1}$、
両側のファイバー次数は3と2となる。どの $q$ にも無限個の三割線は通らない。
実際、$j(q)$ からの平面射影は次数4の基点自由系であり、双有理なら長さ2の分離に失敗する
因子は有限である（各特異像の有限逆像で、重複点は曲線上で一意）。
二重被覆なら像は二次曲線で、そのファイバーの長さ2の因子は全て同じ類を持つため、
$q$ を加えた三割線因子は既に示した一意性に反する。従って両射影は有限で、因子の純性から平坦である。
\end{proof}

\begin{lemma}\label{lem:quintic-cubic-polar}
$V^*$ 上に非退化交代形式 $\omega$ があり、$j(E)$ を切る五次元の三次式系を
$P=(\bigwedge^2 V^*)_0$（$\omega$ による縮約の核）の双対と同一視できる。
曲線上の点 $z$ で三次式が特異となる条件は、対応する汎関数が
$F_z=z\wedge z^\perp$ を消すことである（$z$ は同次座標の代表）。
さらに長さ2の因子 $B$ の張る直線が等方的なら、その直線は三割線である。
\end{lemma}
\begin{proof}
原点を選べば $L=\mathcal O(5O)$ とできる。Weierstrass式 $y^2=x^3+A x+B_0$ に対して
完全系の座標を $(a,b,c,d,e)=(1,x,y,x^2,xy)$ とする。交代五行列の初めの二行の右三成分を
$s=(a,b,c),t=(b,d,e)$、左上の二行二列を零、右下の交代三行列を
ベクトル $(-B_0 a,-d-Aa,c)$ の外積行列とする。
最大Pfaffianは符号を除いて $s,t$ の二次小行列三つと、このベクトルとの内積二つである。
これらの共通零点はスキームとして完全五次曲線である。
$a\ne0$ ではWeierstrass式そのものになり、$a=0$ では $[0:0:0:0:1]$ のみが残る。
$e=1$ で $b=cd,a=bc$ を消去すると、残りは $c-d^2-Aad-B_0 ab=0$ の
滑らかな一方程式となる（他の二式はその倍数）。

中心 $o$ での行列の階数は4である。その核を最後の基底ベクトルに取り、行列を
\[
 \begin{pmatrix} uJ+D(z)&z\\-z^t&0\end{pmatrix}
\]
と書く。最後の列の四切断 $z$ は $o$ からの射影の座標を与える。
独立でなければ $o$ を通る直線全体で同じ列が零となり、そこではPfaffianの条件は一本の二次式だけである。
これは $o$ で非零で、長さ2の割線を与えて矛盾する。$J$ は可逆交代行列である。
初めの四つの符号付きPfaffianは $uJ' z+b(z)$、$J'$ は可逆交代行列
（$J$ のPfaffian余因子）、$b$ は二次の列である。Pfaffian列は元の行列の核に入るので $z^tb(z)=0$。
$v=J'^{-1}b$、$\omega(z,w)=z^tJ'w$ とすると $z\wedge v(z)$ は $P$ に値を取る。
この列の成分は射影曲線をスキームとして切る。実際 $z_i\ne0$ の局所で
$v(z)=-uz$ の $u$ を消去でき、最後のPfaffianも行列の核の恒等式により従う。

各等方二平面 $W$ 上で $v$ の $V^*/W=W^*$ への像は $z$ での評価が零の二次共変ベクトルである。
二変数の一次syzygyにより、これは $z$ に垂直な一次共変ベクトルと一次式 $h_W$ の積である。
$\mathbb P(W)$ が $E$ と交わらなければ $h_W$ の零点で $z\wedge v(z)$ は $W$ を表す。
このような $W$ は一般的である（各 $q\in E$ を通る等方直線は一次元族）。
ゆえに三次写像の像は $\mathbb P(P)$ の滑らかな三次元のPlücker二次曲面に稠密で、五成分は独立である。
全三次系の次元も5である。$j(E)$ は二次曲面に入らない。滑らかな曲面と錐の解消の次数・種数式は
前節のとおりで、次数5では種数1にならない。従って二次の制限は次元10の同型で、
$V$ の基点自由鉛筆の列と $H^1(L)=0$ から三次の制限も全射である。
$z\in E$ では $v(z)$ は $z$ の倍数で、微分像は $F_z$ に入る。
上のスキーム生成と滑らかさにより微分像は二次元、従って $F_z$ 全体となる。
この定数階数の同一視は曲線上の束の同一視である。
最後に $W$ が長さ2の $B$ を含む等方平面なら、$h_W$ は $B$ で零となり恒等零である。
そこで $z\wedge v$ は $W$ 上で一本の三次式に限られ、曲線との交叉は長さ3となる。
\end{proof}

\begin{theorem}\label{thm:elliptic-quintic-centers}
楕円五次点グラフの有限Frobenius候補では $j(E)=1728$（ここで $j$ は楕円曲線の不変量）であり、
標準四移動は平行移動である。各位置で、完全五切断系からの射影中心は、$E$ の四次の単数
$\beta=\pm\mathrm i$ のいずれかについて $\ker(2-\beta)$ の固有直線である。
従って各 $L_i$ を固定すれば中心は高々十点に限られる。
反対二点の階数低下因子の各一般ファイバーは、二点の差が $\ker(1+\mathrm i)$ の
非零元となる被約な対である。
\end{theorem}
\begin{proof}[行列式の三次式が一般に非零であること]
位置0の $p$ を固定し、位置2を変数とした三行三列の行列式を $S_p$ とする。
一般の $p$ で恒等零と仮定する。随伴行列は一般階数1である
（定理\ref{thm:opposite-curve-divisor}により曲線上で元の階数が一般に2）。
四変数多項式環でその二次成分を共通因子と列・行の外積に分解する。
両因子は曲線上で隣接埋込みを一点から射影した平面像を与えるので次数が正で、共に一次となる。
特に $j_1$ を $p$ から射影した平面系と $j_2$ の三つの一次式が同じ写像を与える。
後者には基因子が一点 $x$ だけあり、$L_2(-x)=L_1(-p)$ となるので $x=p+d$、
$d\in E$ は固定である。両平面系は基因子を除けば同じ三切断空間である。

$j_1$ の一般の三割線因子 $T_t$ を取る。補題\ref{lem:elliptic-quintic-trisecants}の有限性により、
その各台 $p$ は上の一般性を満たす。また一般の $T_t$ は少なくとも二つの台を持つ。
全て $3p$ なら標数零でその台は $t$ の有理写像となり、次数9の三倍写像に有理切断を与えてしまう。
$p$ からの平面系では $T_t-p$ が一点へ写る。切断の言葉では、この長さ2の因子で零となる
平面系の部分が二次元以上である。従って $T_t-p+(p+d)$ は $j_2$ で三割線因子である。
$p$ によらず同じ因子類なので、一意性によりこの因子は $p$ によらない。
二つの異なる台で比較すると $d=0$ となる。ゆえに $L_1=L_2$ で、
一般の $p$ で両切断超平面のうち $p$ で消える部分も等しい。
この部分は $p$ とともに動いて各超平面を張るので両系も等しい。
曲線を含む二次曲面がないため、全関係は射影三空間の対角上で零となって矛盾する。

\smallskip
\noindent\textit{階数低下から三割線へ。}
$D_p$ は定理\ref{thm:opposite-curve-divisor}の長さ2の因子とする。
行列の二次小行列が $D_p$ 上で消えるので、$S_p$ はこのスキーム上で特異である。
補題\ref{lem:quintic-cubic-polar}により、これは対応する $P$ 上の汎関数が $F_z$ を $D_p$ 全体で消す条件である。
$D_p$ の張る直線が一般に非等方的だとする。このとき条件は四つ独立で、
汎関数は両点を結ぶ二ベクトルの原始部分に比例する
（外積による非退化な対合で $P$ と $P^*$ を同一視する）。
実際、二点が相異なれば $F_z,F_w$ の共通元は $z\wedge w$ の倍に限られ、非等方なので共通元は零。
二重点なら接方向を $h$ として、$\omega(z,h)\ne0$ である。
$F_z$ の微分を $z\wedge V^*$ で割ると $z\wedge u\mapsto h\wedge u\pmod z$ となり、
$u\in z^\perp/\langle z\rangle$ 上で単射である。この場合の汎関数は $(z\wedge h)_0$ である。
両点の場合の直交恒等式から差を取って対角で制限すればこれが従う。

$p\mapsto [S_p]$ が使う束の次数は高々5である。実際、その係数は三行列式の束
$L_0^3/(L_1L_{-1})=L_0/K$ の切断である。
一方 $p\mapsto \langle j_2(D_p)\rangle$ のPlücker束は
$\det\pi_*(\operatorname{pr}_2^*L_2|_D)$ である（$\pi:D\to E$ は第一射影）。
$D^2=2n$、$n=\Delta\cdot\Gamma_{\sigma^{-1}}$ とおくと、この束の次数は
曲面上の因子列のRiemann--Rochで $10-n$ となる。
$n$ は平行移動なら0、他は $\deg(1-\sigma_*)\le4$ である。
二ベクトルの原始部分への射影中心は階数4の二ベクトルで、Plücker像にないため次数は落ちない。
従ってこの写像の次数は6以上となり、$[S_p]$ と一般に等しいことに反する。
ゆえに一般の $D_p$ は等方直線にあり、補題により三割線に乗る。

\smallskip
\noindent\textit{三割線の入射曲線。}
三割線が切る長さ3の因子の残余一点を $r(p)$、その三次因子類を $t(p)$ とすれば
$p\mapsto (t(p),r(p))$ は $Z$ に値を取る射である。
実際、直線が正則に動き、長さ3の有限交叉は一般に平坦で余因子を取れる。得た有理射を滑らかな基曲線上で延長する。
$t$ は非定数（固定した三割線は有限個の点しか含まない）で、$r$ も補題の有限性で非定数である。
よって $Z$ に、正規化が楕円曲線となる成分がある。
層の重複度も数えれば $Z$ の二つの射影次数は3と2で、各成分は両側に支配的である。
一成分だけなら重複度は $\gcd(3,2)=1$。その正規化からの両射を $f,g$ とすると、
$f$ の三つの逆像の $g$ による像の和が三次因子類に等しい。原点を選び準同型部分を取れば
$f_*=3g_*$（三つの逆像は核元だけ異なる）。これは $3=9\deg g$ を与えて矛盾する。
従って二成分で、どちらも被約かつ第二射影で同型となる（有限双有理射の像の基曲線は正規）。
第一射影の次数は1と2である。原点選択による平行移動定数を略すと、第一成分は
$r=\beta t$、第二成分は $t=\gamma r$ となる。ここで $\beta$ は単数、$\deg\gamma=2$。
三点の和が $t$ なので $1=\beta+\widehat\gamma$ である。二点の和の準同型部分が双対同種写像になるのは
$\widehat\gamma\gamma=[2]$ から従う。
$|1-\beta|^2=2$ より $\beta=\pm\mathrm i$、$j(E)=1728$ である。

第二成分の点を $r(p)$ として除くなら、$D_p$ には同じ成分のもう一点と、第一成分の一点が残る。
これらの $p$ に関する次数はそれぞれ $\deg r$ と $2\deg r$ で、和が第二射影次数2になりえない。
従って除くのは第一成分の点のみで、$D_p$ は $\gamma$ の一ファイバーを与え、
その二点の差は $\ker\gamma=\ker(1+\mathrm i)$ の非零元である。
$D\to E$ の第二射影次数から $\deg(t)=2$、また $t_*-r_*=1+\sigma^{-1}_*$ である。
よって $\deg(1+\sigma^{-1}_*)=\deg(1-\beta)\deg(t)=4$ となり、$\sigma$ の準同型部分は恒等である。

\smallskip
\noindent\textit{射影中心。}
$k$ を $\ker\gamma$ の非零元とすると全三割線因子は
$q+(q+k)+(\beta\gamma q+c)$ と書ける。$\gamma=1-\bar\beta$ なので $\beta\gamma=\beta-1$。
$G=\ker(2-\beta)$ は位数5であり、その各元による平行移動は上の三点を同時に平行移動する。
従ってこの三割線族を保つ。$G$ は次数5の $L$ を保つので完全系の射影四空間に作用する。
各因子の完全系での包は平面であり、これらの平面の共通部分は中心 $o$ 一点である。
他にも共通点があれば射影後に全三割線が同じ点を通る。しかし一般の $q$ を通る相異なる三割線は二本あり、
固定点と $j(q)$ を結ぶ直線の一意性に反する。従って $G$ は $o$ を固定する。
巡回群の生成元を五乗が恒等となるよう線型化すると、$L$ は $E/G$ の次数1の束から引き戻される。
五指標の各固有切断空間はその束の平坦ひねりの切断で一次元なので固有直線は五本である。
位置を回せば全ての中心で同じ結論を得る。
\end{proof}

\begin{corollary}\label{cor:quintic-finite-elliptic-parameters}
$E:y^2=x^3-x$ と固定する。添字を保ち、共通のパラメータ平行移動と各位置の線型座標変更を許すと、
楕円五次窓の全候補は有限個の一パラメータ族の次のリストに含まれる。
四切断系を
\[
 j_\epsilon(x,y)=[1+\epsilon x^2:x:y:xy],\qquad \epsilon=\mathrm i,-\mathrm i
\]
として、位置 $i$ の写像は $p\mapsto j_{\epsilon_i}(p+t_i)$ と取れる。
$t_0=0,t_1=t,t_2=2t+d,t_3=3t+3d+e$、ここで $t\in E,\ d\in E[10],\ e\in E[5]$。
四移動 $p\mapsto p+s$ では $t_4=s$ と選べて $s=4t+6d+f$、$f\in E[5]$ である。
各辺の関係は命題\ref{prop:quintic-mixed-kernel}の全核から復元する。
ここでの有限リストは必要条件で、正則性や全てのパラメータ間の同型は別に決める必要がある。
\end{corollary}
\begin{proof}
$L=\mathcal O(5O)$ の完全座標で $o_\epsilon=[-\epsilon:0:0:1:0]$ とおく。
これは割線外である。実際、$a,b,d$ の部分における有限二点の線型結合で $b=0$ とし、
$c=e=0$ とするには、異なる $x$ の場合は両点の $y=0$ が必要である。
$x=0,1,-1$ の相異なる二つの評価の組からは上の $d/a=\epsilon$ は得られない。
同じ $x$ なら $b=0,a\ne0$ より $x=0$ で $d=0$。一方が $O$ の場合も同様に $b=0$ なら $d=0$。
有限点での接線上では、$o_\epsilon$ は曲線上にないので非零な接方向係数を要する。
$e=xy$ を微分し $c=e=0$ を課すと $y\,dx=0$、ゆえに分岐点である。
そこでは $db=dd=0$ なので $b=0$ は $x=0$ を強制する。$O$ の接線では $a=0$ である。

$k=(0,0)$、$[\epsilon](x,y)=(-x,\epsilon y)$ とする。
$j_\epsilon$ で $q,q+k,([\epsilon]-1)q$ は一般に共線である。実際、座標は
\[
 q+k=(-1/x,y/x^2),\quad
 ([\epsilon]-1)q=\left(\frac{\epsilon(x^2-1)}{2x},\frac{(1-\epsilon)y(x^2+1)}{4x^2}\right)
\]
となり、三行四列の座標の全三次小行列が零である。
三点の和の準同型部分は $1+\epsilon$ で次数2なのでこの族は全ての三割線因子を与える。
前定理の共通平面の議論で $o_\epsilon$ は $\ker(2-\epsilon)$ の固有直線である。
完全系の $E[5]$ 平行移動はこの五つの固有直線を置換し、その作用は推移的である。
実際、$E[5]$ が一つの直線を固定すればその直線に恒等となるよう表現を線型化でき、
$L$ が次数25の商を通じて降下することになり次数5に反する。
位数25の群が五点に固定点なく作用するので軌道は一つである。
一般の $L_i$ は $\mathcal O(5O)$ の平行移動であり、これで二つの系と $t_i$ の表示が従う。

各 $L_i$ の次数5のAbel座標は $-5t_i$ である。平行移動の場合の
系\ref{cor:curve-divisor-monodromy}より二階差 $K$ は二捩れで、添字によらない。
従って $d=t_2-2t_1+t_0$ は十捩れ、次の二階差と $d$ との差は五捩れである。
四番目を含む同じ漸化式を用いれば表示の通りとなる。
\end{proof}

\begin{corollary}\label{cor:quintic-same-center-type}
前系の四つの符号は全て同じで、どの隣接差 $t_{i+1}-t_i$ も五捩れではない。
従って四位置の埋込みは同じ位数5の平行移動群を保つ。
\end{corollary}
\begin{proof}
二つの隣接埋込みが同じ長さ3の三割線因子を持つことはない。
両割線の積 $\mathbb P^1\times\mathbb P^1$ 上で六関係はこの長さ3のグラフに消える。
双一次式の制限はそこへ全射である。実際、どちらの二切断も長さ2を分離する。
三相異点なら二つの一次式の積で任意の二点だけを零にできる。
一つの二重点と単点なら二重点で同時に消える二因子の積で単点を分け、
単点で零かつ二重点で単元の因子と他方の二切断で二重点を分ける。
三重点なら単元と一次の一様化元を両側で選び、積の位数0,1,2を得る。
従って制限された六式は高々一次元を張り、割線の積上に余分な曲線を与える。

各 $j_\epsilon(p+t_i)$ の三割線は $q,q+k,(\epsilon-1)q+c_i$ である。
二符号が違えば、その二つの第三点を一致させる方程式は非零な同種写像の一ファイバーで、必ず解がある。
特殊な $q$ でも因子の一致は重複度つきで成り立つので、これは前段に反する。
従って符号は一定である。

隣接差が五捩れなら両偏極は同じ $L$ である。
中心が同じなら二次曲面がないため $\mathbb P^3$ の対角が点グラフに入る。
中心が異なるなら二中心は同じ巡回群の固有点である。
それらを結ぶ直線が完全系の $E$ に交われば位数5の平行移動の全軌道を含むが、
完全五次系では長さ3の因子は平面を張る。従って二つの四切断空間の共通三切断は基点自由である。
六関係は共通三切断の三交換子を含む。残る三式を完全系の射影四空間に移せば、
共通三座標が非零の開部分で零点は曲線のみで、補座標だけの割線も成分にならない。
割線全体が入れば両側の三座標が共に零になる点対を与え、基点自由性と点グラフの仮定に反する。
従ってこの三二次式の完全交叉の次数8が唯一の曲線成分の次数5の倍数となり矛盾する。
\end{proof}
